\documentclass[12pt]{article}

% Use Times (Times New Roman-like) font
\usepackage{times}

% Set page margins
\usepackage[margin=1in]{geometry}

% For 1.5 line spacing
\usepackage{setspace}

% For custom headers/footers
\usepackage{fancyhdr}

\pagestyle{fancy}
\fancyhf{} % Clear default header and footer
\fancyfoot[C]{\textbf{i}} % Centered page number

% Remove the top horizontal line
\renewcommand{\headrulewidth}{0pt}
\begin{document}
	
	\begin{center}
		{\Large \bf Abstract}
	\end{center}

    \vspace{20pt}
	
	\onehalfspacing % Set 1.5 line spacing for the text that follows
	
	\textit{The rapid advancement of deep learning in medical imaging has led to the development of 
an AI-powered Ultrasound Nerve Imaging and Clinical Decision Support System designed to automate the detection and classification of thyroid disorders, breast cancer, kidney stones, and liver disorders. The system implements a robust three-stage pipeline to enhance diagnostic accuracy and clinical utility. First, preprocessing standardizes input images through resizing and quality enhancement. Next, segmentation using a Convolutional Neural Network precisely isolates regions of interest such as nodules, tumors, or calculi. Finally, disease detection with YOLOv11 classifies abnormalities and generates confidence scores for each prediction. This integrated approach overcomes limitations of traditional methods by combining automated preprocessing with advanced deep learning architectures, eliminating manual intervention and operator dependency. The system delivers comprehensive clinical support by producing diagnostic reports with visual annotations and recommending specialist referrals with contact details and appointment information. Additionally, it provides personalized care guidance including dietary recommendations and educational resources. Experimental results demonstrate superior performance with over 90 percent classification accuracy across all target diseases, significantly outperforming conventional threshold-based techniques. The offline implementation ensures broad accessibility without requiring specialized hardware or real-time processing capabilities. Key innovations include the unified CNN-YOLOv11 architecture for multi-disease detection, end-to-end automation of the diagnostic workflow, and seamless integration of AI analysis with practical healthcare 
recommendations. This research establishes a new standard for AI-assisted ultrasound diagnosis by combining technical precision with clinically actionable outputs, while maintaining scalability for diverse healthcare environments. Future enhancements will focus on expanding disease coverage and optimizing computational efficiency. The system represents a significant step toward standardized, accessible, and accurate ultrasound-based disease detection and management.}
	
	\vspace{1em} % Add vertical space between abstract and keywords
	
	\noindent
	\textbf{Keywords:} Ultrasound imaging, deep learning, disease detection, CNN, YOLOv11, clinical decision support, medical diagnosis.
	
\end{document}
